\newcommand{\coursemodule}{Module 04}
% ════════════════════════════════════════════════════════════════════════════
% Pixels, Frames & Transformers: shared Beamer preamble (WUR style)
% Adapted from the preamble of Bags, Vectors & Transformers.
%
% Use in a module deck (from <part>/<module>/slides/):
%   \newcommand{\coursemodule}{Module 01}
%   % ════════════════════════════════════════════════════════════════════════════
% Pixels, Frames & Transformers: shared Beamer preamble (WUR style)
% Adapted from the preamble of Bags, Vectors & Transformers.
%
% Use in a module deck (from <part>/<module>/slides/):
%   \newcommand{\coursemodule}{Module 01}
%   % ════════════════════════════════════════════════════════════════════════════
% Pixels, Frames & Transformers: shared Beamer preamble (WUR style)
% Adapted from the preamble of Bags, Vectors & Transformers.
%
% Use in a module deck (from <part>/<module>/slides/):
%   \newcommand{\coursemodule}{Module 01}
%   \input{../../../shared/preamble.tex}
%   \begin{document}
%   \moduletitle{Images as Data}
%   ...
% Compile with: latexmk -pdf <deck>.tex
% ════════════════════════════════════════════════════════════════════════════

\documentclass[aspectratio=169]{beamer}

% ── Packages ────────────────────────────────────────────────────────────────
\usepackage[T1]{fontenc}
\usepackage[utf8]{inputenc}
\usepackage{lmodern}
\usepackage{microtype}
\usepackage{xcolor}
\usepackage{tikz}
\usetikzlibrary{arrows.meta, positioning, calc, shapes.geometric, decorations.pathreplacing, matrix, fit}
\usepackage{booktabs}
\usepackage{fontawesome5}
\usepackage{array}
\usepackage[most]{tcolorbox}

\providecommand{\coursemodule}{}

% ── WUR color palette ───────────────────────────────────────────────────────
\definecolor{wurgreen}{HTML}{34B233}
\definecolor{darkgreen}{HTML}{1F6B1E}
\definecolor{lightgreen}{HTML}{E8F7E8}
\definecolor{darknavy}{HTML}{1A1A2E}
\definecolor{midgray}{HTML}{6B7280}
\definecolor{lightgray}{HTML}{F3F4F6}
\definecolor{bodytext}{HTML}{1F2937}
\definecolor{codegreen}{HTML}{9AE6B4}
\definecolor{warnred}{HTML}{FEE2E2}
\definecolor{warnborder}{HTML}{EF4444}

% ── Beamer theme ────────────────────────────────────────────────────────────
\usetheme{default}
\usecolortheme{default}
\usefonttheme{professionalfonts}
\setbeamertemplate{navigation symbols}{}
\setbeamertemplate{itemize items}{\textcolor{wurgreen}{\textbullet}}
\setbeamertemplate{itemize subitem}{\textcolor{midgray}{--}}
\setbeamertemplate{enumerate items}{\textcolor{wurgreen}{\textbf{\insertenumlabel.}}}
\setbeamercolor{frametitle}{fg=white, bg=darknavy}
\setbeamertemplate{frametitle}{%
  \begin{beamercolorbox}[wd=\paperwidth, ht=1.1cm, dp=0.2cm, leftskip=0.6cm]{frametitle}
    \usebeamerfont{frametitle}\insertframetitle
  \end{beamercolorbox}%
}
\setbeamerfont{frametitle}{size=\large, series=\bfseries}
\setbeamercolor{background canvas}{bg=white}
\setbeamercolor{normal text}{fg=bodytext}
\setbeamerfont{normal text}{size=\small}
\setbeamertemplate{footline}{%
  \begin{beamercolorbox}[wd=\paperwidth, ht=0.4cm, dp=0.15cm]{footline}%
    \hspace{0.5cm}%
    \textcolor{midgray}{\tiny Pixels, Frames \& Transformers · \coursemodule\ · Wageningen University \& Research}%
    \hfill%
    \textcolor{midgray}{\tiny \insertframenumber}%
    \hspace{0.4cm}%
  \end{beamercolorbox}%
}

% ── Boxes ───────────────────────────────────────────────────────────────────
\tcbset{wurbase/.style={arc=3pt, boxrule=0.6pt, left=7pt, right=7pt, top=5pt, bottom=5pt,
                        fontupper=\small, before skip=4pt, after skip=4pt}}
\newtcolorbox{greenbox}[1][]{wurbase, colback=lightgreen, colframe=wurgreen, #1}
\newtcolorbox{graybox}[1][]{wurbase, colback=lightgray, colframe=lightgray, #1}
\newtcolorbox{navybox}[1][]{wurbase, colback=darknavy, colframe=darknavy, colupper=white, #1}
\newtcolorbox{warnbox}[1][]{wurbase, colback=warnred, colframe=warnborder, #1}
\newtcolorbox{codebox}[1][]{wurbase, colback=darknavy, colframe=darknavy, colupper=white,
                            fontupper=\ttfamily\scriptsize, left=6pt, top=4pt, bottom=4pt, #1}

% ── Title slide ─────────────────────────────────────────────────────────────
\newcommand{\moduletitle}[1]{%
{
\setbeamertemplate{footline}{}
\begin{frame}[plain]
  \begin{tikzpicture}[remember picture, overlay]
    \fill[darknavy] (current page.south west) rectangle (current page.north east);
    \fill[wurgreen] (current page.south west) rectangle ([xshift=0.45cm]current page.north west);
    \node[anchor=west, text=white, font=\huge\bfseries, text width=13cm]
      at ([xshift=1.1cm, yshift=1.2cm]current page.west) {Pixels, Frames \& Transformers};
    \node[anchor=west, text=wurgreen, font=\large, text width=13cm]
      at ([xshift=1.1cm, yshift=0.05cm]current page.west) {#1 · \coursemodule};
    \node[anchor=west, text=midgray, font=\small, text width=12cm]
      at ([xshift=1.1cm, yshift=-1.5cm]current page.west)
      {Denise J. Roth, PhD Candidate\\
       Strategic Communication Group · Wageningen University \& Research};
  \end{tikzpicture}
\end{frame}
}}

% ── Small helpers ───────────────────────────────────────────────────────────
\newcommand{\keyterm}[1]{\textbf{\textcolor{darkgreen}{#1}}}
\newcommand{\muted}[1]{\textcolor{midgray}{#1}}
\newcommand{\credit}[1]{{\tiny\textcolor{midgray}{#1}}}
\newcommand{\notebookslide}[1]{%
\begin{frame}{In the Notebook}
\vspace{0.1cm}
#1
\end{frame}}
\newcommand{\nbitem}[2]{\item[\textcolor{wurgreen}{#1}] #2}

%   \begin{document}
%   \moduletitle{Images as Data}
%   ...
% Compile with: latexmk -pdf <deck>.tex
% ════════════════════════════════════════════════════════════════════════════

\documentclass[aspectratio=169]{beamer}

% ── Packages ────────────────────────────────────────────────────────────────
\usepackage[T1]{fontenc}
\usepackage[utf8]{inputenc}
\usepackage{lmodern}
\usepackage{microtype}
\usepackage{xcolor}
\usepackage{tikz}
\usetikzlibrary{arrows.meta, positioning, calc, shapes.geometric, decorations.pathreplacing, matrix, fit}
\usepackage{booktabs}
\usepackage{fontawesome5}
\usepackage{array}
\usepackage[most]{tcolorbox}

\providecommand{\coursemodule}{}

% ── WUR color palette ───────────────────────────────────────────────────────
\definecolor{wurgreen}{HTML}{34B233}
\definecolor{darkgreen}{HTML}{1F6B1E}
\definecolor{lightgreen}{HTML}{E8F7E8}
\definecolor{darknavy}{HTML}{1A1A2E}
\definecolor{midgray}{HTML}{6B7280}
\definecolor{lightgray}{HTML}{F3F4F6}
\definecolor{bodytext}{HTML}{1F2937}
\definecolor{codegreen}{HTML}{9AE6B4}
\definecolor{warnred}{HTML}{FEE2E2}
\definecolor{warnborder}{HTML}{EF4444}

% ── Beamer theme ────────────────────────────────────────────────────────────
\usetheme{default}
\usecolortheme{default}
\usefonttheme{professionalfonts}
\setbeamertemplate{navigation symbols}{}
\setbeamertemplate{itemize items}{\textcolor{wurgreen}{\textbullet}}
\setbeamertemplate{itemize subitem}{\textcolor{midgray}{--}}
\setbeamertemplate{enumerate items}{\textcolor{wurgreen}{\textbf{\insertenumlabel.}}}
\setbeamercolor{frametitle}{fg=white, bg=darknavy}
\setbeamertemplate{frametitle}{%
  \begin{beamercolorbox}[wd=\paperwidth, ht=1.1cm, dp=0.2cm, leftskip=0.6cm]{frametitle}
    \usebeamerfont{frametitle}\insertframetitle
  \end{beamercolorbox}%
}
\setbeamerfont{frametitle}{size=\large, series=\bfseries}
\setbeamercolor{background canvas}{bg=white}
\setbeamercolor{normal text}{fg=bodytext}
\setbeamerfont{normal text}{size=\small}
\setbeamertemplate{footline}{%
  \begin{beamercolorbox}[wd=\paperwidth, ht=0.4cm, dp=0.15cm]{footline}%
    \hspace{0.5cm}%
    \textcolor{midgray}{\tiny Pixels, Frames \& Transformers · \coursemodule\ · Wageningen University \& Research}%
    \hfill%
    \textcolor{midgray}{\tiny \insertframenumber}%
    \hspace{0.4cm}%
  \end{beamercolorbox}%
}

% ── Boxes ───────────────────────────────────────────────────────────────────
\tcbset{wurbase/.style={arc=3pt, boxrule=0.6pt, left=7pt, right=7pt, top=5pt, bottom=5pt,
                        fontupper=\small, before skip=4pt, after skip=4pt}}
\newtcolorbox{greenbox}[1][]{wurbase, colback=lightgreen, colframe=wurgreen, #1}
\newtcolorbox{graybox}[1][]{wurbase, colback=lightgray, colframe=lightgray, #1}
\newtcolorbox{navybox}[1][]{wurbase, colback=darknavy, colframe=darknavy, colupper=white, #1}
\newtcolorbox{warnbox}[1][]{wurbase, colback=warnred, colframe=warnborder, #1}
\newtcolorbox{codebox}[1][]{wurbase, colback=darknavy, colframe=darknavy, colupper=white,
                            fontupper=\ttfamily\scriptsize, left=6pt, top=4pt, bottom=4pt, #1}

% ── Title slide ─────────────────────────────────────────────────────────────
\newcommand{\moduletitle}[1]{%
{
\setbeamertemplate{footline}{}
\begin{frame}[plain]
  \begin{tikzpicture}[remember picture, overlay]
    \fill[darknavy] (current page.south west) rectangle (current page.north east);
    \fill[wurgreen] (current page.south west) rectangle ([xshift=0.45cm]current page.north west);
    \node[anchor=west, text=white, font=\huge\bfseries, text width=13cm]
      at ([xshift=1.1cm, yshift=1.2cm]current page.west) {Pixels, Frames \& Transformers};
    \node[anchor=west, text=wurgreen, font=\large, text width=13cm]
      at ([xshift=1.1cm, yshift=0.05cm]current page.west) {#1 · \coursemodule};
    \node[anchor=west, text=midgray, font=\small, text width=12cm]
      at ([xshift=1.1cm, yshift=-1.5cm]current page.west)
      {Denise J. Roth, PhD Candidate\\
       Strategic Communication Group · Wageningen University \& Research};
  \end{tikzpicture}
\end{frame}
}}

% ── Small helpers ───────────────────────────────────────────────────────────
\newcommand{\keyterm}[1]{\textbf{\textcolor{darkgreen}{#1}}}
\newcommand{\muted}[1]{\textcolor{midgray}{#1}}
\newcommand{\credit}[1]{{\tiny\textcolor{midgray}{#1}}}
\newcommand{\notebookslide}[1]{%
\begin{frame}{In the Notebook}
\vspace{0.1cm}
#1
\end{frame}}
\newcommand{\nbitem}[2]{\item[\textcolor{wurgreen}{#1}] #2}

%   \begin{document}
%   \moduletitle{Images as Data}
%   ...
% Compile with: latexmk -pdf <deck>.tex
% ════════════════════════════════════════════════════════════════════════════

\documentclass[aspectratio=169]{beamer}

% ── Packages ────────────────────────────────────────────────────────────────
\usepackage[T1]{fontenc}
\usepackage[utf8]{inputenc}
\usepackage{lmodern}
\usepackage{microtype}
\usepackage{xcolor}
\usepackage{tikz}
\usetikzlibrary{arrows.meta, positioning, calc, shapes.geometric, decorations.pathreplacing, matrix, fit}
\usepackage{booktabs}
\usepackage{fontawesome5}
\usepackage{array}
\usepackage[most]{tcolorbox}

\providecommand{\coursemodule}{}

% ── WUR color palette ───────────────────────────────────────────────────────
\definecolor{wurgreen}{HTML}{34B233}
\definecolor{darkgreen}{HTML}{1F6B1E}
\definecolor{lightgreen}{HTML}{E8F7E8}
\definecolor{darknavy}{HTML}{1A1A2E}
\definecolor{midgray}{HTML}{6B7280}
\definecolor{lightgray}{HTML}{F3F4F6}
\definecolor{bodytext}{HTML}{1F2937}
\definecolor{codegreen}{HTML}{9AE6B4}
\definecolor{warnred}{HTML}{FEE2E2}
\definecolor{warnborder}{HTML}{EF4444}

% ── Beamer theme ────────────────────────────────────────────────────────────
\usetheme{default}
\usecolortheme{default}
\usefonttheme{professionalfonts}
\setbeamertemplate{navigation symbols}{}
\setbeamertemplate{itemize items}{\textcolor{wurgreen}{\textbullet}}
\setbeamertemplate{itemize subitem}{\textcolor{midgray}{--}}
\setbeamertemplate{enumerate items}{\textcolor{wurgreen}{\textbf{\insertenumlabel.}}}
\setbeamercolor{frametitle}{fg=white, bg=darknavy}
\setbeamertemplate{frametitle}{%
  \begin{beamercolorbox}[wd=\paperwidth, ht=1.1cm, dp=0.2cm, leftskip=0.6cm]{frametitle}
    \usebeamerfont{frametitle}\insertframetitle
  \end{beamercolorbox}%
}
\setbeamerfont{frametitle}{size=\large, series=\bfseries}
\setbeamercolor{background canvas}{bg=white}
\setbeamercolor{normal text}{fg=bodytext}
\setbeamerfont{normal text}{size=\small}
\setbeamertemplate{footline}{%
  \begin{beamercolorbox}[wd=\paperwidth, ht=0.4cm, dp=0.15cm]{footline}%
    \hspace{0.5cm}%
    \textcolor{midgray}{\tiny Pixels, Frames \& Transformers · \coursemodule\ · Wageningen University \& Research}%
    \hfill%
    \textcolor{midgray}{\tiny \insertframenumber}%
    \hspace{0.4cm}%
  \end{beamercolorbox}%
}

% ── Boxes ───────────────────────────────────────────────────────────────────
\tcbset{wurbase/.style={arc=3pt, boxrule=0.6pt, left=7pt, right=7pt, top=5pt, bottom=5pt,
                        fontupper=\small, before skip=4pt, after skip=4pt}}
\newtcolorbox{greenbox}[1][]{wurbase, colback=lightgreen, colframe=wurgreen, #1}
\newtcolorbox{graybox}[1][]{wurbase, colback=lightgray, colframe=lightgray, #1}
\newtcolorbox{navybox}[1][]{wurbase, colback=darknavy, colframe=darknavy, colupper=white, #1}
\newtcolorbox{warnbox}[1][]{wurbase, colback=warnred, colframe=warnborder, #1}
\newtcolorbox{codebox}[1][]{wurbase, colback=darknavy, colframe=darknavy, colupper=white,
                            fontupper=\ttfamily\scriptsize, left=6pt, top=4pt, bottom=4pt, #1}

% ── Title slide ─────────────────────────────────────────────────────────────
\newcommand{\moduletitle}[1]{%
{
\setbeamertemplate{footline}{}
\begin{frame}[plain]
  \begin{tikzpicture}[remember picture, overlay]
    \fill[darknavy] (current page.south west) rectangle (current page.north east);
    \fill[wurgreen] (current page.south west) rectangle ([xshift=0.45cm]current page.north west);
    \node[anchor=west, text=white, font=\huge\bfseries, text width=13cm]
      at ([xshift=1.1cm, yshift=1.2cm]current page.west) {Pixels, Frames \& Transformers};
    \node[anchor=west, text=wurgreen, font=\large, text width=13cm]
      at ([xshift=1.1cm, yshift=0.05cm]current page.west) {#1 · \coursemodule};
    \node[anchor=west, text=midgray, font=\small, text width=12cm]
      at ([xshift=1.1cm, yshift=-1.5cm]current page.west)
      {Denise J. Roth, PhD Candidate\\
       Strategic Communication Group · Wageningen University \& Research};
  \end{tikzpicture}
\end{frame}
}}

% ── Small helpers ───────────────────────────────────────────────────────────
\newcommand{\keyterm}[1]{\textbf{\textcolor{darkgreen}{#1}}}
\newcommand{\muted}[1]{\textcolor{midgray}{#1}}
\newcommand{\credit}[1]{{\tiny\textcolor{midgray}{#1}}}
\newcommand{\notebookslide}[1]{%
\begin{frame}{In the Notebook}
\vspace{0.1cm}
#1
\end{frame}}
\newcommand{\nbitem}[2]{\item[\textcolor{wurgreen}{#1}] #2}


\begin{document}

\moduletitle{Coding Images and Classic Classifiers}

% ─────────────────────────────────────────────────────────────────────────────
\begin{frame}{Human Coding Is the Benchmark}
Every automated measure in this course, from a brightness score to a vision-language model, is
only as good as its agreement with careful human judgment.

\vspace{0.3cm}
\centering
\begin{tikzpicture}[
  b/.style={draw=wurgreen, fill=lightgreen, rounded corners=3pt, font=\footnotesize, align=center,
            inner sep=4pt, text width=1.95cm, minimum height=1.1cm},
  n/.style={draw=darknavy, fill=darknavy, text=white, rounded corners=3pt, font=\footnotesize\bfseries,
            align=center, inner sep=4pt, text width=1.95cm, minimum height=1.1cm},
  arr/.style={-{Stealth[length=2mm]}, thick, color=midgray}]
\node[b] (cb) {codebook};
\node[b, right=0.5cm of cb] (pilot) {pilot coding};
\node[b, right=0.5cm of pilot] (rel) {reliability good enough?};
\node[n, right=0.75cm of rel] (gold) {coded test set};
\node[b, right=0.5cm of gold] (auto) {score any automated method};
\draw[arr] (cb) -- (pilot); \draw[arr] (pilot) -- (rel);
\draw[arr] (rel) -- node[above, font=\scriptsize, color=midgray] {yes} (gold);
\draw[arr] (gold) -- (auto);
\draw[arr] (rel.south) -- ++(0,-0.45) -| node[pos=0.25, below, font=\scriptsize, color=midgray]
      {no: revise the codebook} (cb.south);
\end{tikzpicture}

\vspace{0.4cm}
\raggedright
\begin{greenbox}
In this module you code 40 images. Those codes become the \textbf{test set} for every method in
Parts 1 to 3, so all of them are scored on the same images.
\end{greenbox}
\end{frame}

% ─────────────────────────────────────────────────────────────────────────────
\begin{frame}{Choosing the Unit of Analysis}
The same picture can be coded at several levels. Each level answers different questions.

\vspace{0.2cm}
\begin{columns}[T]
\begin{column}{0.48\textwidth}
\begin{greenbox}
\textbf{\faImage\ The whole image}\\
\emph{Does this photo show a protest?}
\end{greenbox}
\begin{greenbox}
\textbf{\faUser\ Each person}\\
\emph{How are women shown in news photos?}
\end{greenbox}
\end{column}
\begin{column}{0.48\textwidth}
\begin{greenbox}
\textbf{\faCrop\ A region}\\
\emph{What is in the foreground? How much of the package is the mascot?}
\end{greenbox}
\begin{greenbox}
\textbf{\faComments\ The post}\\
\emph{Does the caption change what the image means?}
\end{greenbox}
\end{column}
\end{columns}

\vspace{0.1cm}
\muted{This module codes whole images. Regions and people return in modules 12 and 14, posts in
module 11.}
\end{frame}

% ─────────────────────────────────────────────────────────────────────────────
\begin{frame}{A Short Visual Codebook}
\vspace{0.1cm}
{\footnotesize
\begin{tabular}{@{}>{\bfseries}l>{\raggedright\arraybackslash}p{6.1cm}>{\raggedright\arraybackslash}p{5.2cm}@{}}
\toprule
 & \textbf{``yes'' when \ldots} & \textbf{``no'' when \ldots} \\
\midrule
protest & people engaged in protest are visible: marching, gathering with signs, a blockade, a sit-in & only objects (a poster, an empty square), police alone, or a crowd with no sign of protest \\[2pt]
crowd & roughly ten or more people are visible, including the background & fewer than about ten people \\[2pt]
police & officers, police vehicles or police barriers are visible & no police (stewards don't count) \\[2pt]
signs & a sign, banner or flag with a message is carried or displayed by people & no signs, or only street, shop and traffic signs \\
\bottomrule
\end{tabular}}

\vspace{0.3cm}
\begin{navybox}
\textbf{Code what you can see, not what you know about the event.} If you can't tell, choose the
answer that needs less inference, and write down the image: it probably calls for a new rule.
\end{navybox}
\end{frame}

% ─────────────────────────────────────────────────────────────────────────────
\begin{frame}{Coding in Practice}
\begin{columns}[T]
\begin{column}{0.47\textwidth}
\textbf{Tools}

\vspace{0.1cm}
\begin{graybox}
\textbf{In the notebook}: buttons under each image, codes saved after every click. Quick, for
whole-image codes and one coder.
\end{graybox}
\begin{graybox}
\textbf{Label Studio} (open source): boxes and regions, several coders, a web interface. Worth
setting up for a real project.
\end{graybox}
\begin{graybox}
\textbf{A spreadsheet} with file names: the fallback that always works.
\end{graybox}
\end{column}
\begin{column}{0.5\textwidth}
\textbf{Habits}
\begin{itemize}
\item \keyterm{Code blind}: don't let coders see where an image came from
\item \keyterm{Shuffle} the order, and code in short sessions
\item \keyterm{Write down new rules} as hard cases come up
\item \keyterm{Recode} a sample later: is your coding \emph{stable}?
\item Add a \keyterm{second coder}: is it \emph{reproducible}? This is the standard for
      publication.
\end{itemize}
\end{column}
\end{columns}
\end{frame}

% ─────────────────────────────────────────────────────────────────────────────
\begin{frame}{Measuring Agreement}
\begin{columns}[T]
\begin{column}{0.5\textwidth}
\begin{itemize}
\item \keyterm{Percent agreement}: share of images coded the same way; rewards luck
\item \keyterm{Cohen's kappa}: corrected for chance,
      $\kappa = (p_o - p_e)\,/\,(1 - p_e)$
\item \keyterm{Krippendorff's alpha}: any number of coders, missing codes; the standard in
      content analysis
\end{itemize}
\begin{graybox}
$\alpha \geq .800$: reliable \quad $.667$--$.800$: tentative conclusions only
\muted{(Krippendorff, 2018)}
\end{graybox}
\end{column}
\begin{column}{0.46\textwidth}
\vspace{0.3cm}
\begin{warnbox}
\textbf{Why percent agreement misleads}\\
\emph{Police} appears in 5\% of images. One coder looks carefully; the other always says ``no''
without looking.\\[4pt]
\begin{tabular}{@{}lr@{}}
percent agreement & \textbf{0.95} \\
Cohen's kappa & \textbf{0.00} \\
\end{tabular}
\end{warnbox}
\muted{Report a chance-corrected measure. Rare categories need more images before you can judge
their reliability at all.}
\end{column}
\end{columns}
\end{frame}

% ─────────────────────────────────────────────────────────────────────────────
\begin{frame}{Free Labels Are Noisy Labels}
Every image comes with a label you didn't code: the \textbf{category} it was filed under.
Metadata, hashtags, search terms and accounts work the same way.

\vspace{0.25cm}
\centering
\begin{tikzpicture}[
  cat/.style={draw=darknavy, fill=darknavy, text=white, rounded corners=3pt, font=\footnotesize,
              align=center, inner sep=5pt, text width=3.4cm},
  ok/.style={draw=wurgreen, fill=lightgreen, rounded corners=3pt, font=\footnotesize, align=left,
             inner sep=4pt, text width=4.2cm},
  bad/.style={draw=warnborder, fill=warnred, rounded corners=3pt, font=\footnotesize, align=left,
             inner sep=4pt, text width=4.2cm},
  arr/.style={-{Stealth[length=2mm]}, thick, color=midgray}]
\node[cat] (c) {Category:\\ \textbf{Demonstrations and protests in the Netherlands in 2019}};
\node[ok, right=1.4cm of c, yshift=1.1cm] (a) {\faCheck\ a march with banners};
\node[bad, below=0.12cm of a] (b) {\faTimes\ a poster on a wall};
\node[bad, below=0.12cm of b] (d) {\faTimes\ a police van on an empty square};
\node[bad, below=0.12cm of d] (e) {\textbf{?}\ a speaker alone on a stage};
\foreach \n in {a,b,d,e} \draw[arr] (c.east) -- (\n.west);
\end{tikzpicture}

\vspace{0.3cm}
\raggedright
\begin{greenbox}
A category says \textbf{where an image was filed}, not \textbf{what it shows}. Train on free
labels if nothing better exists, but always evaluate against human codes.
\end{greenbox}
\end{frame}

% ─────────────────────────────────────────────────────────────────────────────
\begin{frame}{A Classic Classifier, Honestly Evaluated}
\begin{columns}[T]
\begin{column}{0.5\textwidth}
\centering
\begin{tikzpicture}[
  tr/.style={draw=wurgreen, fill=lightgreen, rounded corners=3pt, font=\footnotesize, align=center,
             inner sep=5pt, text width=5.4cm},
  te/.style={draw=darknavy, fill=darknavy, text=white, rounded corners=3pt, font=\footnotesize,
             align=center, inner sep=5pt, text width=5.4cm},
  arr/.style={-{Stealth[length=2mm]}, thick, color=midgray}]
\node[tr] (train) {\textbf{All other images}\\ category labels \\ 5-fold cross-validation \\ fit the final model};
\node[te, below=0.5cm of train] (test) {\textbf{Your 40 coded images}\\ never seen in training\\ scored against your codes};
\draw[arr] (train) -- node[right, font=\scriptsize, color=midgray] {predict} (test);
\end{tikzpicture}
\end{column}
\begin{column}{0.47\textwidth}
\begin{itemize}
\item Input: the 11 features from module 03
\item \keyterm{Logistic regression}: one weight per feature
\item \keyterm{Random forest}: many trees; thresholds and interactions
\item Unequal classes: balanced class weights; report \keyterm{balanced accuracy} (0.5 = chance)
      and \keyterm{F1} per class
\end{itemize}
\end{column}
\end{columns}

\vspace{0.2cm}
\begin{warnbox}
A model scored on images it was trained on looks far better than it is. Keep the test set out of
training and out of every modeling decision.
\end{warnbox}
\end{frame}

% ─────────────────────────────────────────────────────────────────────────────
\begin{frame}{The Semantic Gap}
\centering
\begin{tikzpicture}[
  l/.style={draw=midgray, fill=lightgray, rounded corners=3pt, font=\footnotesize, align=left,
            inner sep=6pt, text width=4.6cm},
  r/.style={draw=darknavy, fill=darknavy, text=white, rounded corners=3pt, font=\footnotesize,
            align=left, inner sep=6pt, text width=4.6cm}]
\node[l] (look) {\textbf{How it looks}\\[2pt] light, color, edges, symmetry,\\ roughly how many frontal faces};
\node[r, right=3.2cm of look] (show) {\textbf{What it shows}\\[2pt] people protesting, signs with\\ messages, police, a celebration};
\draw[very thick, dashed, color=warnborder] (look.east) -- node[above, font=\scriptsize, color=warnborder, text width=2.8cm, align=center]
     {only through correlations}
     node[below, font=\scriptsize, color=midgray, text width=2.8cm, align=center] {orange $\approx$ King's Day} (show.west);
\end{tikzpicture}

\vspace{0.4cm}
\raggedright
Handcrafted features can only cross this gap where meaning happens to correlate with the look,
and those correlations break on new data: an orange-clad protest becomes ``King's Day''.

\vspace{0.2cm}
\begin{greenbox}
\textbf{Next, Part 2:} features that represent content. Networks that learn their own features
(module 05) and features borrowed from networks trained on millions of images (module 06), scored
on the same 40 images.
\end{greenbox}
\end{frame}

% ─────────────────────────────────────────────────────────────────────────────
\notebookslide{
\begin{itemize}
\setlength{\itemsep}{5pt}
\nbitem{\faClipboardList}{A four-variable codebook, and the hard cases it needs rules for}
\nbitem{\faTags}{Coding 40 images with a labeling tool in the notebook}
\nbitem{\faRedo}{Recoding 20 of them: stability with kappa and alpha}
\nbitem{\faFolderOpen}{Category labels against your own codes}
\nbitem{\faLaptopCode}{Logistic regression and a random forest on handcrafted features}
\nbitem{\faSearch}{Testing against your codes and reading the mistakes}
\end{itemize}

\vspace{0.35cm}
\muted{\footnotesize part1\_pixels/04\_coding\_images\_and\_classic\_classifiers/notebooks/\\
04\_coding\_images\_and\_classic\_classifiers\_exercises.ipynb · 6 exercises · CPU runtime ·
about 20 minutes of coding}
}

\end{document}
